\chapter*{บรรณานุกรม}
\addcontentsline{toc}{chapter}{บรรณานุกรม}

\begin{list}{}{%
  \setlength{\leftmargin}{1.25cm}%
  \setlength{\itemindent}{-1.25cm}%
  \setlength{\parsep}{\parskip}%
  \setlength{\itemsep}{8pt plus 2pt minus 1pt}%
}

\item American College of Sports Medicine. (2021). \textit{ACSM's guidelines for exercise testing and prescription} (11th ed.). Wolters Kluwer.

\item Fogg, B. J. (2009). A behavior model for persuasive design. \textit{Proceedings of the 4th International Conference on Persuasive Technology} (Persuasive '09), Article 40, 1--7. \url{https://doi.org/10.1145/1541948.1541999}

\item Mifflin, M. D., St Jeor, S. T., Hill, L. A., Scott, B. J., Daugherty, S. A., \& Koh, Y. O. (1990). A new predictive equation for resting energy expenditure in healthy individuals. \textit{The American Journal of Clinical Nutrition}, \textit{51}(2), 241--247. \url{https://doi.org/10.1093/ajcn/51.2.241}

\item Thassana, C. (2023). Physics-informed edge intelligence and sensor integration for digital health and precision kinematics. \textit{Journal of Applied Physics and Sensor Technology}, \textit{15}(3), 112--128.

\item Thassana, C., \& Janthamalee, J. (2024). Dynamic thresholding and refractory windowing in MEMS accelerometer gait analysis. \textit{IEEE Sensors Journal Letters}, \textit{8}(2), 45--52.

\item Tudor-Locke, C., Craig, C. L., Brown, W. J., Clemes, S. A., De Cocker, K., Giles-Corti, B., Hatano, Y., Inoue, S., Matsudo, S. M., Mutrie, N., Oppert, J. M., Rowe, D. A., Schmidt, M. D., Schofield, G. M., Spence, J. C., Teixeira, P. J., Tully, M. A., \& Blair, S. N. (2011). How many steps/day are enough? for adults. \textit{International Journal of Behavioral Nutrition and Physical Activity}, \textit{8}(1), 79. \url{https://doi.org/10.1186/1479-5868-8-79}

\item World Health Organization. (2020). \textit{WHO guidelines on physical activity and sedentary behaviour}. World Health Organization.

\item ชีวะ ทัศนา. (2567). \textit{ฟิสิกส์เชิงคำนวณและระบบตรวจวัดสัญญาณอัจฉริยะสำหรับวิทยาการข้อมูล}. สำนักพิมพ์มหาวิทยาลัยราชภัฏรำไพพรรณี.

\item ชีวะ ทัศนา, และ คณะนักวิจัย. (2569). \textit{การออกแบบและพัฒนาแอปพลิเคชันเพื่อสุขภาพ การนับก้าวและการตรวจจับการเคลื่อนไหวร่างกายเชิงสรีรวิทยาประยุกต์}. หน่วยวิเคราะห์-ติดตามคุณภาพดินและสภาพอากาศเชิงพื้นที่ คณะวิทยาศาสตร์และเทคโนโลยี มหาวิทยาลัยราชภัฏรำไพพรรณี.

\end{list}
